\begin{figure}[H]
    \centering
    \begin{tikzpicture}[
        node distance=0.8cm and 1cm,
        stepbox/.style={rectangle, draw, rounded corners, minimum width=6cm, minimum height=1.2cm, align=center, thick, fill=orange!10},
        arrow/.style={-stealth, thick},
        font=\small
    ]

    % Draw Grid for GPU
    \draw[thick, draw=green!60!black, rounded corners, loosely dashed] (-4.5, 1.5) rectangle (4.5, -6.5);
    \node[anchor=north west, font=\bfseries\large, text=green!50!black] at (-4.2, 1.2) {GPU (CUDA) 并行架构化};

    % Phase 1: Line Graph Construction
    \node[stepbox] at (0, -0.5) (csr) {第一阶段: 拓扑高并发构建 \\ (CSR度数扫描与前缀和无锁映射)};
    
    % Phase 2: Conditional Weight SIMT
    \node[stepbox] at (0, -3.0) (simt) {第二阶段: 条件权重提取 \\ (SIMT多分支协同与并发访存优化)};
    
    % Phase 3: $\Delta$-Stepping
    \node[stepbox] at (0, .5.5) (delta) {第三阶段: $\Delta$-Stepping寻优 \\ (无序波前松弛与共享内存并发探底)};

    \draw[arrow] (csr) -- (simt);
    \draw[arrow] (simt) -- (delta);

    \end{tikzpicture}
    \caption{基于CUDA架构的线图寻路并行优化策略实施框架}
    \label{fig:parallel_optimization}
\end{figure}